\section{Three Deployment Decisions}\label{sec:discussion}
The three checks in Figure~\ref{fig:defense-decision} answer concrete design decisions. Table~\ref{tab:synthesis-decisions} collects the resulting permissions and boundaries so that each case can focus on what changes in the deployment judgment.

\subsection{Case 1: choose the final interface to protect}
\paragraph{Decision and relevant evidence}
An application designer must decide whether upstream API padding or batching also protects the user-facing connection after a content-aware frontend. NetEcho identifies the relevant receiving path: Trace A is token counts per group; Trace B estimates token/group character lengths. A downstream application can consume decoded upstream content or stream state and emit finer-grained network observations \cite{netecho2025}. Captured deployment configurations and locally implemented active-defense settings are distinct source results. Neither exact text nor original generation timestamps are automatically recovered.

\paragraph{What changes after binding}
An upstream padded link answers a different question from protection of the final emitter. NetEcho's authors already establish re-emission in their systems. The joint synthesis specifies the next measurement: capture the last target-bearing interface rather than merely refit on the upstream trace. UI Traces offers a contrasting case: modifying stored timestamps preserves action types and order, so the first question is which surviving features predict identity \cite{uitraces2026}. Refitting that view cannot certify a different downstream one.

SplitDecoy supplies an independent receiving-path countercondition: decoys obscure a forward split-model view, yet zero/nonzero support in the returned gradient identifies the real rows on every tested frame across nine seeds \cite[Sections 5.1--5.3]{politis2026splitdecoy}. Its authors already identify this defect. The synthesis places it beside NetEcho: inspect every target-bearing outgoing path, not only the interface on which protection was evaluated. Section~\ref{sec:controlled-placement} independently demonstrates loss and restoration of size equalization at a new terminal emitter.

\paragraph{When transfer is permitted}
The source-local inference result applies when the receiver obtains the tested projection under its evaluated population and condition. A stronger information-theoretic permission has a precise form.

\noindent\textbf{Proposition 1 (protected-only postprocessing).}
Let $S$ be the hidden target, $Z$ the complete protected representation and $Y$ the downstream observation. If $S\to Z\to Y$ is a Markov chain, including any randomization independent of $S$ given $Z$, then
\[
 I(S;Y)\leq I(S;Z).
\]
This is the data processing inequality \cite[Theorem~2.8.1]{cover2006elements}: expand $I(S;Z,Y)$ in both orders, use $I(S;Y\mid Z)=0$ and non-negativity of conditional mutual information. Thus a verified protected-only transformation preserves an applicable mutual-information upper bound. A frontend with additional decoded content or target-dependent state need not satisfy that premise; NetEcho illustrates this second path. The proposition licenses an information bound, not transfer of a particular classifier's accuracy.

\subsection{Case 2: choose whether cover may execute}
\paragraph{Decision and relevant evidence}
WRA provides a privacy-oriented reason to execute cover: create plausible alternative research activity around a user's real request. It issues real and decoy prompts in parallel, returns only the real output, and evaluates interleaved visited-domain sets. With one decoy, its functional-intent reconstruction metric falls from 0.770 to 0.610 \cite{jeong2026wra}. This supports reducing the tested observer's prompt inference. AHB supplies a complementary task-state result: Trip Planning success is 0.75 Raw versus 0.15 for combined Online+Fake humanization, with unintended GUI transitions discussed by the authors \cite{zhu2026turing}. The studies measure different endpoints: WRA tests private-intent inference; AHB tests human--agent distinction and task performance.

\paragraph{What changes after binding}
Four execution classes make the choice concrete: (i) transport dummies discarded before dispatch, (ii) requests confined to an isolated shadow environment, (iii) read-only/idempotent operations whose initial effects must be checked, and (iv) state-changing actions. Section~\ref{sec:controlled-placement} implements these routes, using a pure read for the third class. Idempotence constrains repetition; it does not make the first state change safe. WRA establishes a privacy benefit for its domain-view construction, while AHB establishes a state risk for one combined humanization condition. Together they motivate separating \emph{whether cover obscures the target} from \emph{whether its execution preserves the task}. WRA's separate visibility-condition usefulness scores are not paired terminal utility for its decoy experiment.

\paragraph{When transfer is permitted}
Retain WRA's inference benefit for its tested domain view and decoy construction. Verified discard or non-effecting isolation excludes direct cover-induced task changes; execution against the real environment instead calls for objective goal and policy checks. Same-condition terminal utility and numerical latency/wire cost remain \NRR{} for WRA's decoy result. Dangol's live loopback OHTTP/A2A condition separately provides a protocol budget: accuracy 0.30 at about 13$\times$ bytes and 11$\times$ latency \cite{dangol2026workflow}. AHB's non-Fake History+Long point (task success 0.7143; detection 0.9905) further shows that state effects and indistinguishability depend on the configuration. These local findings guide which cover implementation to test, rather than pricing one implementation with another's results.

\subsection{Case 3: choose a valid adaptation result and price}
\paragraph{Decision and relevant evidence}
A service owner choosing cache protection under a workload budget needs both a valid protected-condition attack estimate and a cost for the actual matching mechanism. An adaptation result from another receiving view answers only its own subproblem. UI Traces compares clean- and delay-trained classifiers on the same delayed archive view. That permits the positive source-local conclusion that adaptation recovers substantial prediction under the tested transformation (D2) \cite{uitraces2026}. SIGLEAK's useful procedural inference under a common SkillGuard5 prompt is a different result, not an installation effect or a paired adaptation contrast \cite{geng2026sigleak}.

Early Bird's $K$-specific hit/miss calibration and subsequent KV-oracle simulation are distinct from its GPTCache de-identification prototype's 6ms/4\% overhead \cite{earlybird2025}. KVGov measures a 1.3/1.6$\mu$s median/p99 hash-based message authentication code (HMAC) increment, while workload-dependent lost cache reuse remains unmeasured; key-equality isolation does not establish isolation of a semantic nearest-neighbor index \cite{addagada2026kvgov}. SpliceDefense likewise separates its constant-time fusion signal effect from its quantized-padding cost \cite{spliceleak2026}. For full-prompt recovery, Early Bird Table~VI also distinguishes queries per recovered token from recovery coverage and all-token cost; the authors already explain this survivor-conditioned denominator.

\paragraph{What changes, and when transfer is permitted}
Allow the UI adaptation contrast for its matched target, receiver, transformation and defended test units. Allow each component cost for the measured component and mechanism. CachePrune provides a positive cache-protection case: marked-token recovery, task quality and serving performance are evaluated for selective sharing, while contextual inference and missed sensitive tokens remain separate conditions \cite{wu2026cacheprune}. OTRO permits a similarly precise claim: epoch padding reduces access-count information \emph{beyond byte length} from 3.94 to 0.04 bits in its 20K-prompt closed world, not total prompt information \cite{lee2026otro}. A new receiver, workload, combined mechanism or matching primitive needs a corresponding bridge. The added output is the choice of which result may inform which decision, including these positive local uses, rather than a ranking of unlike milliseconds or microseconds.

The supplementary archive reanalysis supplies a compact illustration of the same rule: refitting after the specified offline timestamp perturbation improves paired loss by 0.1503 nats on the same defended records (Supplement S4). It illustrates adaptation within a fixed receiving view; the source cases supply the intervention-placement evidence.

\subsection{Evidence strength and the resulting decision rule}
Table~\ref{tab:synthesis-decisions} distinguishes direct support, extensions and background. NetEcho anchors re-emission; SplitDecoy tests a different unprotected return path. WRA anchors privacy-oriented decoy execution, and AHB anchors the mobile state-effect contrast. The controlled experiments add direct local tests of the first two decisions. UI Traces, Early Bird, KVGov and SpliceDefense anchor adaptation/component distinctions; CachePrune and OTRO supply positive mechanism-bound uses. Each conclusion is tied to its measurement. Supplements S6/S10 report source-deletion and expansion sensitivity.
\begin{table*}[t]\centering\footnotesize
\caption{Deployment judgments, positive reuse conditions and applicable boundaries. Direct evidence tests the named branch; extensions address a different mechanism, mode or component.}\label{tab:synthesis-decisions}
\begin{tabularx}{\textwidth}{@{}p{.16\textwidth}p{.27\textwidth}YY@{}}\toprule
Decision & Direct support / extension / background & Judgment after binding & Positive license and remaining bridge\\\midrule
Protect the final stream & Direct: NetEcho IV-C--E/Fig.5; our paired two-link test. Extension: SplitDecoy 5.1--5.3, backward support; UI Traces 6.1, surviving actions. & An upstream-view result does not certify a new content-aware emitter. & Keep the result for the tested view. Verify the same projection, or evaluate the final hop; protected-only postprocessing preserves an applicable information bound.\\
Implement cover as actions & Direct: WRA 8.1/Table8, private intent; AHB 5.2.3/Table1/E.3, mobile state effects; our four-route test. Extension: Dangol 9.7 live transport costs. & Decoys have measured privacy value; execution safety is a separate task-state question. & Keep WRA's domain-view inference result. Verified discard/shadow/pure reads preserve state in our test; live cover requires goal/policy and resource checks.\\
Choose adaptation and cost & Direct: UI Traces 6.1, matched adaptation. Positive local cases: CachePrune 5; OTRO VII. Extensions: Early Bird, KVGov, SpliceDefense component/mode evidence. & Permit matched adaptation, selective-cache results and conditional tokenizer reduction; distinguish their measured cost objects. & Match training/test access, mechanism and workload. OTRO byte length remains visible; its epoch-padding estimate is separate from main TTFT figures.\\\bottomrule
\end{tabularx}\end{table*}


The practical payoff is choosing the right next measurement: final-hop capture, task-state validation, or mechanism/workload evaluation. The following inventory makes these decisions explicit across result records.
